% !TEX root = ../risk_vs.tex
\section{Data-Mined Predictability}\label{sec:dm}
We describe our data mining procedure and the predictability it uncovers.

\subsection{Data Mining Procedure}\label{sec:dm-data}

We begin with 241 Compustat annual accounting variables used in \citet{yan2017fundamental}. Yan and Zheng select these variables to (1) ensure non-missing values in at least 20 years and (2) that the average number of firms with non-missing values is at least 1,000 per year. We add CRSP market equity, leading to 242 ``ingredient'' variables. A more sophisticated selection would filter on data availability in real time. But given that data availability changes in large, positive, and permanent jumps (\citealt{easterwood2024do}), the more complicated procedure would likely yield similar results.

We then generate 29,315 accounting ratios (signals) using two functional forms: simple ratios ($X/Y$) and first differences scaled by a lagged denominator ($\Delta X / \text{lag}(Y)$). The numerator can use any of the 242 ingredients. The denominator is restricted to the 65 ingredients that are not zero for at least 25\% of firms in 1963 with matched CRSP data. This procedure leads to $\approx 242 \times 65 \times 2 = 31,460$ ratios, but we drop 2,145 ratios that are redundant in ``unsigned'' portfolio sorts.\footnote{For the $65\times65 = 4,225$ ratios where the numerator is also a valid denominator, there are only 65 choose 2 = 2,080 ratios that are in a sense distinct.}  

We lag each signal by six months relative to the fiscal year end, and then form long-short decile strategies by sorting stocks on the lagged signals in each June.  Delisting returns and other data handling methods follow \citet{ChenZimmermann2021}.  For further details, please see \url{https://github.com/chenandrewy/flex-mining}.

% motivations
This procedure aims to be the simplest possible data mining benchmark for peer-reviewed predictors. Accounting data is the modal data source for peer-reviewed predictors, representing roughly 50\% of the CZ dataset. The second most common data source is past returns, which represents only 20\%.  While it is possible to data mine both accounting data and past returns simultaneously (e.g. \citealt{chen2023high}), it requires several additional design choices and significantly complicates the analysis. 

Alternatively, one can motivate this procedure as a data mining benchmark that could have been constructed \emph{before} the bulk of the modern literature. The idea that accounting ratios may be informative about firm value goes back to the 1930s: 17 of the 26 chapters on stock selection in \citet{graham1934securityanalysis} focus on accounting statements. \citet{smith1935changes} examine 21 accounting ratios for their ability to predict financial distress (see also \citealt{ramser1931}; \citealt{fitzpatrick1932}; and \citealt{merwin1942}). These works were available decades before \citet{fama1973risk}. By the 50th anniversary of \citet{graham1934securityanalysis}, the book had become quite influential, and was celebrated in Warren \citepos{buffett1984superinvestors} popular speech and article. Thus, one well could have decided to search accounting ratios in 1984, when only 9 of the 212 predictors in the CZ dataset had been published. Indeed, \citet{ou1989financial} propose selecting stocks based on ``a large number of financial statement attributes,'' using a process in which ``[n]o conscious attempt is made to assess predictive ability on the basis of what we think should work.''

Our selection of accounting ratios contrasts with Yan and Zheng \citeyearpar{yan2017fundamental}, who use functional forms inspired, in part, by the asset pricing literature. Our procedure avoids the concern that such inspiration induces look-ahead bias. Nevertheless, previous versions of this paper used Yan and Zheng's data and found similar results. 

\subsection{Out-of-Sample Returns from Data Mining}\label{sec:dm-datasum}

Our data mining procedure generates notable out-of-sample returns, as seen in Panel (a) of Table \ref{tab:dm-oos}.  Each June, we sort the 29,000 accounting ratios into five bins based on their mean long-short returns over the past 30 years (in-sample) and compute the mean return over the next year within each bin (out-of-sample).  We then average these statistics across each year.  

\begin{table}[!h]
\caption{Descriptive Statistics of Data-Mined Accounting Strategies}
\label{tab:dm-oos}

\begin{singlespace}
\noindent The table describes data-mined accounting strategies using ``out-of-sample'' portfolio sorts (Panel (a)) and PCA (Panel (b)). Panel (a) sorts all ratios each June into 5 bins based on past 30-year long-short returns (in-sample) and computes the mean return over the next year within each bin (out-of-sample). Statistics are calculated by strategy, then averaged within bins, then averaged across sorting years. Decay is the percentage decrease in mean return out-of-sample relative to in-sample. We omit decay for bin 4 because the mean return in-sample is negligible.  Panel (b) applies PCA to ratios that have t-statistics greater than 2.0 in at least 10\% of the in-sample periods from Panel (a). Data-mined predictors resemble published ones in terms of in-sample performance, out-of-sample performance, and covariance structure.
\end{singlespace}
\begin{centering}
\vspace{-1ex}
\par\end{centering}
\centering{}\setlength{\tabcolsep}{0.9ex} \small
\begin{center}
% ==== begin paste
% Table generated by Excel2LaTeX from sheet 'dm summary'
\begin{tabular}{crrrrrrrrrrrr} \toprule 
\multicolumn{13}{c}{Panel (a): ``Out-of-Sample'' Returns of All Ratios 1994-2020} \\ 
\midrule
In- &   & \multicolumn{5}{c}{Equal-Weighted Long-Short Deciles} &   & \multicolumn{5}{c}{Value-Weighted Long-Short Deciles} 
\\ \cmidrule{3-7}\cmidrule{9-13}Sample &   & \multicolumn{2}{c}{Past 30 Years (IS)} &   & \multicolumn{2}{c}{Next Year (OOS)} &   & \multicolumn{2}{c}{Past 30 Years (IS)} &   & \multicolumn{2}{c}{Next Year (OOS)} \\ 
\cmidrule{3-4}\cmidrule{6-7}\cmidrule{9-10}\cmidrule{12-13}Bin &   & \multicolumn{1}{c}{Return} & \multicolumn{1}{c}{\multirow{2}[2]{*}{t-stat}} &   & \multicolumn{1}{c}{Return} & \multicolumn{1}{c}{Decay} &   & \multicolumn{1}{c}{Return} & \multicolumn{1}{c}{\multirow{2}[2]{*}{t-stat}} &   & \multicolumn{1}{c}{Return} & \multicolumn{1}{c}{Decay} \\ 
&   & \multicolumn{1}{c}{(bps pm)} &   &   & \multicolumn{1}{c}{(bps pm)} & \multicolumn{1}{c}{(\%)} &   & \multicolumn{1}{c}{(bps pm)} &   &   & \multicolumn{1}{c}{(bps pm)} & \multicolumn{1}{c}{(\%)} \\ 
\cmidrule{1-1}\cmidrule{3-4}\cmidrule{6-7}\cmidrule{9-10}\cmidrule{12-13}
% latex table generated in R 4.5.3 by xtable 1.8-4 package
% Thu Aug 20 09:55:48 2026
 1 &  & -59.5 & -4.09 &  & -41.6 & 30.1 &  & -37.6 & -1.96 &  & -15.9 & 57.7 \\ 
  2 &  & -28.5 & -2.35 &  & -16.7 & 41.4 &  & -15.3 & -0.97 &  & -4.7 & 68.9 \\ 
  3 &  & -12.6 & -1.10 &  & -4.3 & 65.9 &  & -4.7 & -0.31 &  & -3.6 & 23.6 \\ 
  4 &  & 0.7 & 0.05 &  & 4.1 &  &  & 5.1 & 0.34 &  & -1.8 &  \\ 
  5 &  & 22.8 & 1.44 &  & 14.5 & 36.4 &  & 25.2 & 1.30 &  & 4.7 & 81.6 \\ 
  
\bottomrule \end{tabular}% 
\vspace{2ex}

% pca panel
\setlength{\tabcolsep}{1.5ex}
% latex table generated in R 4.5.3 by xtable 1.8-4 package
% Thu Aug 20 10:14:59 2026
\begin{tabular}{clccccccccccc}
  \toprule
 & \multicolumn{11}{c}{Panel (b): PCA Explained Variance (\%)} \\
  \midrule
Number of PCs & 1 & 5 & 10 & 20 & 30 & 40 & 50 & 60 & 70 & 80 & 90 & 100 \\ 
\midrule
  Equal-Weighted & 23 & 50 & 58 & 67 & 72 & 75 & 78 & 81 & 83 & 84 & 86 & 87 \\ 
  Value-Weighted & 16 & 41 & 50 & 61 & 68 & 73 & 77 & 80 & 82 & 85 & 87 & 88 \\ 
   \bottomrule
\end{tabular}


\end{center} 
\end{table}

Using equal-weighted strategies, the in-sample returns of the first bin are on average -59 bps per month, with an average t-stat of -4.2.  These statistics are similar to those of the typical published predictor (\citet{ChenZimmermann2021}).  Out-of-sample, the first bin returns -47 bps per month, implying a  decay of only  20\%, once again resembling published predictability (\citet{mclean2016does}).  Since investors can flip the long and short legs of these strategies, these statistics imply substantial out-of-sample returns.  Similar predictability is seen in bin 5, which decays by 32\%.  

Out-of-sample predictability is also seen in value-weighted strategies but with smaller magnitudes. Still, the out-of-sample returns monotonically increase in the in-sample return, indicating the presence of true predictability.  Moreover, the roughly  60\% decay is far from 100\%, and is in the ballpark of the post-sample decay for published predictors.  Similarly, out-of-sample predictability is much weaker post-2004, though it still exists (see Appendix Table \ref{tab:dm-sum-2004}). The concentration of predictability in small stocks and the pre-2004 sample is also found in published predictors (\citealt{chen2022zeroing}). 

\subsection{Data-Mined Predictability Themes}\label{sec:dm-themes}

Since there are 29,000 data-mined signals, Panel (a) of Table \ref{tab:dm-oos} implies thousands of strategies with notable out-of-sample predictability. But how many distinct themes are in these strategies? 

Panel (b) of Table \ref{tab:dm-oos} helps address this question. It applies PCA to the predictive ratios: ratios that have t-statistics greater than 2.0 in at least 10\% of the 30-year in-sample periods from Panel (a). 

PCA results in a non-trivial factor structure: the first five PCs explain about 50\% of total variance. However, many dozens of PCs are required to span 80\% of total variance. A similar variance decomposition is seen in published predictors (\citealt{kozak2018interpreting}; \citealt{bessembinder2021time}; \citealt{chen2023missing}). Pairwise correlations lead to a similar conclusion (Internet Appendix Table \ref{tab:dm-cor}).

% named themes methods
One can alternatively examine the themes by examining the numerators that generate the very largest t-statistics. Table \ref{tab:dm-theme} does this by reporting the 20 numerator and stock weight (equal- or value-) combinations that produce the largest mean t-stats, where the mean is taken across the 65 possible denominators. The table uses the 1963-1980 sample, but similar themes are found in other samples (Internet Appendix \ref{sec:intapp-themes}).

\begin{table}
\caption{Themes from Mining Accounting Ratios in 1980}\label{tab:dm-theme}
\small

Table reports the 20 accounting ratio numerator and stock weight (equal- or value-) combinations with the largest mean t-stats using returns in the years 1963-1980 (IS). `ew' is equal-weight, `vw' is value-weight. We manually group numerators into themes from the literature. Strategies are signed to have positive mean returns IS. `Pct Short' is the share of strategies that short stocks with high  ratios.  `t-stat' and  `Mean Return' are averages across the 65 possible denominators. `Mean Return' is in bps per month. `Mean return OOS/IS' is the mean in either 1981-2004 or 2005-2022 (OOS),  divided by the mean IS. Data mining can uncover themes from the literature before they are published.
\vspace{2ex}

\centering

\setlength{\tabcolsep}{2.0ex}

\begin{tabular}{lrrrlrr}
\toprule
 & \multicolumn{3}{c}{1963-1980 (IS)} &  & \multicolumn{1}{c}{1981-2004} & \multicolumn{1}{c}{2005-2023} \\ \cmidrule{2-4} \cmidrule{6-7}
Numerator (Stock Weight) & \multicolumn{1}{c}{Pct} & \multirow{2}{*}{t-stat} & \multicolumn{1}{c}{Mean} &  & \multicolumn{2}{c}{Mean Return} \\  
 & \multicolumn{1}{c}{Short} &  & \multicolumn{1}{c}{Return} &  & \multicolumn{2}{c}{OOS / IS} \\ 
\midrule
\multicolumn{7}{l}{Investment / Investment Growth (Titman, Wei, Xie 2004; Cooper, Gulen, Schill 2008)} \\ \hline
$\Delta$PPE net (ew) & 98 & 3.9 & 0.80 &  & 1.07 & 0.19\\
$\Delta$Assets (ew) & 100 & 3.9 & 0.87 &  & 1.04 & 0.32\\
$\Delta$Intangible assets (ew) & 100 & 3.8 & 0.52 &  & 1.04 & 0.26\\
$\Delta$PPE gross (ew) & 98 & 3.7 & 0.78 &  & 0.98 & 0.14\\
$\Delta$Invested capital (ew) & 98 & 3.3 & 0.71 &  & 1.36 & 0.35\\
$\Delta$Capital expenditure (ew) & 100 & 3.1 & 0.43 &  & 1.54 & 0.46 \bigstrut[b] \\ 
\multicolumn{7}{l}{External Financing (Spiess and Affleck-Graves 1999; Pontiff and Woodgate 2008)} \\ \hline
$\Delta$Common stock (ew) & 100 & 5.1 & 0.81 &  & 0.66 & 0.34\\
$\Delta$Liabilities (ew) & 100 & 4.6 & 0.81 &  & 0.78 & 0.28\\
$\Delta$Capital surplus (ew) & 100 & 4.2 & 0.61 &  & 1.19 & 0.99\\
$\Delta$Long-term debt (ew) & 100 & 3.4 & 0.45 &  & 1.50 & 0.24\\
$\Delta$Capital surplus (vw) & 98 & 3.0 & 0.54 &  & 0.93 & 0.54 \bigstrut[b] \\ 
\multicolumn{7}{l}{Accruals / Inventory Growth (Sloan 1996; Thomas and Zhang 2002; Belo and Lin 2012)} \\ \hline
$\Delta$Inventories (ew) & 100 & 4.0 & 0.66 &  & 1.22 & 0.22\\
$\Delta$Notes payable st (ew) & 100 & 3.8 & 0.45 &  & 0.56 & 0.25\\
$\Delta$Debt in current liab (ew) & 100 & 3.8 & 0.45 &  & 0.70 & 0.27\\
$\Delta$Current liabilities (ew) & 100 & 3.7 & 0.52 &  & 1.30 & 0.21\\
$\Delta$Receivables (ew) & 100 & 3.5 & 0.65 &  & 0.61 & 0.34 \bigstrut[b] \\ 
\multicolumn{7}{l}{Earnings Surprise (Foster, Olsen, Shevlin 1984; Chan, Jegadeesh, Lakonishok 1996)} \\ \hline
$\Delta$Cost of goods sold (ew) & 100 & 3.5 & 0.60 &  & 0.87 & 0.23\\
$\Delta$Operating expenses (ew) & 100 & 3.5 & 0.60 &  & 0.96 & 0.34\\
$\Delta$SG\&A (ew) & 100 & 3.5 & 0.67 &  & 0.96 & 0.24\\
$\Delta$Interest expense (ew) & 98 & 3.2 & 0.47 &  & 1.38 & 0.73\\
\bottomrule
\end{tabular}

\end{table}

% describe themes
All of the top 20 numerators fit into themes from the cross-sectional literature.  These themes include investment (\citet{titman2004capital}), debt issuance (\citet{spiess1999long}), share issuance (\citet{loughran1995new}), accruals (\citet{sloan1996stock}),  inventory growth (\citet{thomas2002inventory}), and earnings surprise (\citet{watts1978systematic}). For all of these themes, the sign of predictability obtained from  data mining is the same as the sign from the literature (e.g. short stocks with high investment).  Notably, most of these themes are published after 1980, sometimes long after.

% discuss decay
The predictive power of these themes persists out-of-sample (``OOS/IS'' columns). All data-mined numerators produce positive mean returns after 1980. In fact, the return decay in the 1981-2004 out-of-sample period is on average zero. 

% hint at the main results
Taken together, these results hint at our main finding. Data-mined predictability resembles that of peer-reviewed research. This resemblance is seen in performance both in- and out-of-sample (Table \ref{tab:dm-oos}, Panel (a)), covariance structure (Table \ref{tab:dm-oos}, Panel (b)), and themes (Table \ref{tab:dm-theme}). 


